% =============================================================================
% 01-motivation.tex  —  Section: Κίνητρο & Στόχος  (~3 frames, ~3 min)
% =============================================================================
\section{Κίνητρο \& Στόχος}

% --- Frame S1: quantum threat ---
\begin{frame}{Οι κβαντικοί υπολογιστές ακυρώνουν τη γέννηση κρυπτογραφία δημόσιου κλειδιού}
  \begin{columns}[T]
    \begin{column}{0.56\textwidth}
      \begin{itemize}
        \item \textbf{Αλγόριθμος Shor (1994):} RSA, ECC σπάνε σε \en{polynomial time}
        \item \defterm{\en{Harvest-now, decrypt-later}}: τα δεδομένα συλλέγονται \emph{σήμερα}
        \item \textbf{\en{NIST PQC}} στανταρδοποίηση: \en{ML-KEM, ML-DSA, SLH-DSA} (2024) +\\
              \HL{\en{HQC} επιλέχθηκε ως 4ο \en{KEM}} (Μάρτιος 2025)
        \item \HQC: code-based, \en{lattice}-\emph{ανεξάρτητη} ασφάλεια, \en{diversity} αντί \en{ML-KEM}
        \item Σύνθημα βάρους: \en{code/compute} \emph{πολύ μεγαλύτερο} από \en{ML-KEM}
      \end{itemize}
    \end{column}
    \begin{column}{0.42\textwidth}
      \centering
      \resizebox{\columnwidth}{!}{%
        % Timeline: Classical Cryptography → Quantum Threat → NIST Competition → Standardization
% Used in §1.1 to visualize the urgency of PQC migration
\tikzsetnextfilename{ch1-introduction-timeline-pqc}
% Okabe-Ito colour-blind-safe palette (review 1.1 #9)
\definecolor{oiOrange}{HTML}{E69F00}
\definecolor{oiGreen}{HTML}{009E73}
\definecolor{oiBlue}{HTML}{0072B2}
\definecolor{oiVermillion}{HTML}{D55E00}
\begin{tikzpicture}[
  x=2.8cm,
  every node/.style={font=\small},
  era/.style={fill=#1!15, draw=#1!60, rounded corners=2pt,
              minimum height=1.2cm, text width=2.4cm, align=center,
              font=\small\sffamily},
  milestone/.style={circle, fill=#1!80, inner sep=2pt, minimum size=6pt},
  yearlab/.style={font=\footnotesize\sffamily, below=3pt},
  annotation/.style={font=\scriptsize, text width=2.2cm, align=center},
]

% Main timeline axis
\draw[thick, -{Stealth[length=5pt]}] (-0.3,0) -- (5.3,0);

% Era blocks (above the line)
\node[era=gray]   at (0.5, 1.2) {Κλασική\\κρυπτογραφία\\(RSA, ECC)};
\node[era=oiVermillion]    at (1.8, 1.2) {Κβαντική\\απειλή\\(Shor, 1994)};
\node[era=oiOrange] at (3.1, 1.2) {NIST PQC\\Competition\\(2016–2024)};
\node[era=oiGreen] at (4.4, 1.2) {Τυποποίηση\\PQC\\(2024–2035)};

% Milestone dots on the axis
\node[milestone=gray]          (m1) at (0.5, 0) {};
\node[milestone=oiVermillion]           (m2) at (1.8, 0) {};
\node[milestone=oiOrange]        (m3) at (3.1, 0) {};
\node[milestone=oiGreen](m4) at (4.4, 0) {};

% Year labels below
\node[yearlab] at (m1) {1978};
\node[yearlab] at (m2) {1994};
\node[yearlab] at (m3) {2016};
\node[yearlab] at (m4) {2024};

% Additional deadline marker
\node[milestone=oiGreen] (m5) at (5.0, 0) {};
\node[yearlab] at (m5) {2035};

% Annotations below the axis
\node[annotation, below=14pt] at (m1) {RSA\\δημοσιεύεται};
\node[annotation, below=14pt] at (m2) {Αλγόριθμος\\Shor};
\node[annotation, below=14pt] at (m3) {NIST Call\\for Proposals};
\node[annotation, below=14pt] at (m4) {ML-KEM/DSA\\τυποποιούνται};
\node[annotation, below=14pt] at (m5) {CNSA 2.0\\deadline};

% Connecting arrows from eras to milestones
\foreach \m in {m1,m2,m3,m4} {
  \draw[gray!50, thin] (\m) -- ++(0, 0.55);
}

% HNDL threat annotation (spanning arrow)
\draw[oiVermillion, thick, decorate, decoration={brace, amplitude=4pt, mirror}]
  (1.8, -1.6) -- (5.0, -1.6)
  node[midway, below=6pt, font=\scriptsize\sffamily, text=oiVermillion, text width=3cm, align=center]
  {Harvest Now, Decrypt Later\\(δεδομένα ήδη σε κίνδυνο)};

\end{tikzpicture}
%
      }
    \end{column}
  \end{columns}
  \pnote{60 s. Σταθερή: Shor threat already real, NIST picked HQC as backup to lattice diversity.}
\end{frame}

% --- Frame S2: RISC-V opportunity + spectrum ---
\begin{frame}{Επιτάχυνση \en{HQC} \emph{χωρίς} ειδική hardware: οι επικυρωμένες επεκτάσεις \en{ISA} του \en{RISC-V}}
  \begin{columns}[T]
    \begin{column}{0.50\textwidth}
      \begin{itemize}
        \item \defterm{Φάσμα επιτάχυνσης}: \en{C} $\to$ \en{ISA ext} $\to$ \en{FPGA} $\to$ \en{ASIC}
        \item \en{RISC-V}: ανοιχτή, modular \en{ISA} με \en{cryptographic agility}
        \item \HL{\Zbc/\Zvbc}: \en{carry-less} πολ\textperiodcenteredπολ. ($\GFpoly$) — \emph{κρίσιμος πυρήνας}
        \item \Zbb/\Zvbb: \en{bit-manipulation} \& \en{rotate} (για \en{Keccak})
        \item \HL{Μηδέν custom \en{hardware}} — καθαρό λογισμικό, επικυρωμένες επεκτάσεις ISA
      \end{itemize}
    \end{column}
    \begin{column}{0.48\textwidth}
      \centering
      \resizebox{\columnwidth}{!}{%
        % VARIANT — not \input by main.tex; promote manually.
% _v1 — beautified original + TOKEN MIGRATION of acceleration-spectrum.tex.
% Same 4-box flexibility<->performance layout; inline fonts/weights migrated to
% thesis-tikz-styles.tex tokens, the "you are here" pill uses the figpill token,
% and the focal box uses figbox/focal weights. FIX: the original used an
% UNDEFINED `figaccentleader' style for the pill connector (no such token) —
% replaced with `draw=figaccent, figLW/base'. Kept inside \resizebox so the wide
% layout never overflows the text block.
\resizebox{\textwidth}{!}{%
\tikzsetnextfilename{ch1-introduction-acceleration-spectrum_v1}
\begin{tikzpicture}[
  node distance=0.8cm,
  box/.style={
    draw=figrule, rounded corners=3pt, minimum width=3cm, minimum height=1.4cm,
    align=center, figT/label, fill=white
  },
  highlight/.style={box, draw=figaccent, figLW/emph, fill=figaccentfill},
  endlabel/.style={figT/annot},
  rowtitle/.style={font=\sffamily\footnotesize\bfseries, text=figneutral, anchor=west},
  ch/.style={figT/annot, align=center},
]

% --- Boxes (left = most flexible, right = least flexible) ---
\node[box] (sw) {Γενικός κώδικας C\\[-2pt]\footnotesize (base ISA)};
\node[highlight, right=0.8cm of sw] (isa) {ISA Extensions\\[-2pt]\footnotesize (Zbc, Zvbc, Zvbb)};
\node[box, right=0.8cm of isa] (fpga) {FPGA\\[-2pt]\footnotesize Accelerator};
\node[box, right=0.8cm of fpga] (asic) {Custom ASIC\\[-2pt]\footnotesize (σταθερή λειτουργία)};

% --- Gradient spectrum bar (flexibility[green] -> performance[vermillion]) ---
\shade[left color=figsolution!35, right color=figthreat!42, rounded corners=2pt]
  ([yshift=0.70cm]sw.north west) rectangle ([yshift=1.18cm]asic.north east);
\node[endlabel, anchor=west] at ([xshift=2pt,yshift=0.94cm]sw.north west)
  {$\leftarrow$ Μεγαλύτερη ευελιξία};
\node[endlabel, anchor=east] at ([xshift=-2pt,yshift=0.94cm]asic.north east)
  {Μεγαλύτερη απόδοση $\rightarrow$};

% --- "You are here" pill on the ISA box (figpill token) ---
\node[figpill, below=0.55cm of isa] (here) {Παρούσα εργασία};
\draw[draw=figaccent, figLW/base] (isa.south) -- (here.north);

% --- "Χρόνος αλλαγής": row heading + per-box values ---
\coordinate (titley) at ([yshift=-1.9cm]sw.south);
\node[rowtitle] at (sw.west |- titley) {Χρόνος αλλαγής:};
\coordinate (rowy) at ([yshift=-2.5cm]sw.south);
\node[ch] at (sw |- rowy)   {λεπτά–ώρες};
\node[ch] at (isa |- rowy)  {λεπτά (recompile)};
\node[ch] at (fpga |- rowy) {ημέρες–εβδομάδες};
\node[ch] at (asic |- rowy) {μήνες–έτη};

\end{tikzpicture}%
}% end resizebox
%
      }
    \end{column}
  \end{columns}
  \pnote{60 s. Spectrum: our sweet spot = ISA ext, no silicon tapeout. Rationale for RISC-V modularity.}
\end{frame}

% --- Frame S3: contribution statement ---
\begin{frame}{Συνεισφορά: η πρώτη βελτιστοποιημένη \en{pure-software} υλοποίηση \HQC σε \en{RISC-V}}
  \begin{columns}[T]
    \begin{column}{0.50\textwidth}
      \begin{itemize}
        \item \textbf{3-\en{tier} στρατηγική}: \en{scalar Zbc} $\to$ \en{Base-RVV} $\to$ \en{Zvbc}
        \item Πολ\textperiodcenteredπολ.: \HL{\speedup{264}} (\en{Tier~3, HQC-128})
        \item Πλήρης \KEM: \HL{\speedup{76}} (\en{HQC-128}), \speedup{164} (\en{HQC-256})
        \item \en{Hardware}: \speedup{46}/\speedup{78} επικυρωμένα σε \en{BPI-F3} (\en{SpacemiT K1})
        \item Ανοιχτός κώδικας: \code{github.com/fraidakis/hqc-riscv-thesis}
      \end{itemize}
    \end{column}
    \begin{column}{0.48\textwidth}
      \centering
      \resizebox{\columnwidth}{!}{%
        % VARIANT (_v1: baseline refined — token migration + canonical tier palette)
% — not \input by main.tex; promote manually. Same layout as the original.
% Fig 6.1 (sec:polymul-tiers): hardware-tier superset diagram.
\tikzsetnextfilename{ch6-polymul-polymul-tiers_v1}
\begin{tikzpicture}[
  tier/.style={
    draw, rounded corners=4pt, minimum width=11cm, minimum height=1.6cm,
    align=center, figT/title, figLW/base
  },
  % local: mono at the micro scale (no mono+scriptsize token in the set)
  extlabel/.style={font=\ttfamily\scriptsize, anchor=west},
  % local: figbrace token + mirror (token has no mirrored variant)
  lbrace/.style={decorate, decoration={brace, amplitude=6pt, raise=2pt, mirror},
                 draw=figrule, figLW/base},
]

% Tier 1 (top)
\node[tier, fill=figtier1!12, draw=figtier1!70!black] (t1) at (0, 4.4)
  {Tier 1 — Scalar\\[2pt]{\normalfont\ttfamily\scriptsize rv64gc + Zbc/Zbb}};
% Tier 2 (middle)
\node[tier, fill=figtier2!12, draw=figtier2!70!black] (t2) at (0, 2.2)
  {Tier 2 — Base-RVV\\[2pt]{\normalfont\ttfamily\scriptsize rv64gcv + Zbc/Zbb}};
% Tier 3 (bottom)
\node[tier, fill=figtier3!12, draw=figtier3!70!black] (t3) at (0, 0)
  {Tier 3 — Full Vector-Crypto\\[2pt]{\normalfont\ttfamily\scriptsize rv64gcv + Zbc/Zbb/Zvbc/Zvbb}};

% Extension annotations (right)
\node[extlabel] at (5.9, 4.4) {clmul, clmulh, rori, andn};
\node[extlabel] at (5.9, 2.2) {+ vle64, vse64, vxor, vsetvl};
\node[extlabel] at (5.9, 0)   {+ vclmul.vx, vclmulh.vx, vror.vv};

% Variant labels (left) — colour-matched to canonical tier palette
\node[figT/label, anchor=east, text=figtier1!70!black] at (-5.9, 4.4) {\texttt{scalar\_zb}};
\node[figT/label, anchor=east, text=figtier2!70!black] at (-5.9, 2.2) {\texttt{vector\_rvv}};
\node[figT/label, anchor=east, text=figtier3!70!black] at (-5.9, 0)   {\texttt{vector\_zvk}};

% Superset arrows
\draw[figedge] (t1.south) -- (t2.north)
  node[midway, right, figT/annot] {$+$ RVV 1.0};
\draw[figedge] (t2.south) -- (t3.north)
  node[midway, right, figT/annot] {$+$ Zvbc/Zvbb};

\end{tikzpicture}
%
      }
    \end{column}
  \end{columns}
  \pnote{60 s. Land the headline numbers clearly. Committee: three tiers, all pure-software.}
\end{frame}
